% Part: computability
% Chapter: computability-theory
% Section: complement-ce

\documentclass[../../../include/open-logic-section]{subfiles}

\begin{document}

\olfileid{cmp}{thy}{cmp}
\olsection{په محاسبوي ډول د شمېر وړ سټونه د متمم لاندې تړلي نۀ دي}

فرض کړئ $A$ په محاسبوي ډول د شمېر وړ دے۔ ايا د~$A$ متمم،
$\Complement{A} = \Nat \setminus A$، همېشه په محاسبوي ډول د شمېر وړ
هم دے؟ لاندې قضيه او نتيجه ښيي چې ځواب ``نه'' دے۔

\begin{thm}
\ollabel{thm:ce-comp}
فرض کړئ $A$ د طبيعي عددونو هر يو سټ دے۔ بيا $A$ هله او يوازې هله
محاسبه کېدونکے دے چې $A$ او~$\Complement{A}$ دواړه په محاسبوي ډول
د شمېر وړ وي۔
\end{thm}

\begin{proof}
مخکښ لورے اسان دے: که~$A$ محاسبه کېدونکے وي، نو~$\Complement{A}$
هم محاسبه کېدونکے دے ($\Char{A} = 1 \tsub
\Char{\Complement{A}}$)، او ځکه دواړه په محاسبوي ډول د شمېر وړ دي۔

په بل لوري کښې فرض کړئ~$A$ او~$\Complement{A}$ دواړه په محاسبوي ډول
د شمېر وړ دي۔ $A$ د~$\cfind{d}$ د تعريف ساحه، او
$\Complement{A}$ د~$\cfind{e}$ د تعريف ساحه وي۔ تابع~$h$ داسې تعريف
کړئ:
\[
h(x) = \umin{s}{(T(d,x,s) \lor T(e,x,s))}.
\]
په بله وينا، پر ننوت~$x$ تابع~$h$ يا د~$\cfind{d}$ يا
د~$\cfind{e}$ يوه درېدلې محاسبه لټوي۔ اوس که $x \in A$ وي، په
لومړي حالت کښې به بريالۍ شي، او که $x \in \Complement{A}$ وي، په
دويم حالت کښې به بريالۍ شي۔ نو~$h$ هرځاے تعريف شوې محاسبه کېدونکې
تابع ده۔ خو اوس د هر~$x$ دپاره $x \in A$ هله او يوازې هله چې
$T(d, x, h(x))$ وي، يعنې که~$\cfind{d}$ پر همدې ننوت تعريف شوې وي۔
څرنګه چې $T(d, x, h(x))$ محاسبه کېدونکې اړيکه ده، نو~$A$
محاسبه کېدونکے دے۔
\end{proof}

\begin{explain}
په غير رسمي محاسبوي ژبه کښې خبره لا اسانه پوهېدل کېږي: د~$A$ د
پرېکړې دپاره پر ننوت~$x$ د~$\cfind{d}$ او~$\cfind{e}$ محاسبې په
نوبت، يوه يوه ګام، پر مخ بوځئ۔ يوه به خامخا ودرېږي؛ که~$\cfind{d}$ ودرېږي، نو~$x$
په~$A$ کښې دے، او که نه، $x$ په~$\Complement{A}$ کښې دے۔
\end{explain}

\begin{cor}
\ollabel{cor:comp-k}
$\Complement{K_0}$ په محاسبوي ډول د شمېر وړ نۀ دے۔
\end{cor}

\begin{proof}
موږ پوهېږو چې~$K_0$ په محاسبوي ډول د شمېر وړ دے، خو محاسبه کېدونکے
نۀ دے۔ که~$\Complement{K_0}$ په محاسبوي ډول د شمېر وړ وای، نو د
\olref{thm:ce-comp} له مخې به~$K_0$ محاسبه کېدونکے و، چې له
\olref[ncp]{thm:K-0} سره تناقض دے۔
\end{proof}

\end{document}
